\section{確率と統計的推論}
\label{b11_estimation}

\subsection{授業内容}
\subsubsection{概要}
母集団に確率分布を仮定し，そこから得られる標本分布の特徴を利用して母数を推定する方法について学ぶ。標本の平均値が母平均に一致することを確認し，また標本の分散は母分散に一致しないことを確認する。これまでは手元にあるデータを記述する方法を学んできたが，本コマからは未知の母集団の特性を推測するための統計的推論の基礎に入る。現実の研究では全数調査は困難であるため，標本から母集団へと推論することが必要となる。ここでの推定方法はモーメントを用いたものであるが，確率分布を用いる方法もあり，その場合の``確率モデル''の考え方が必要なことを学ぶ。

\subsubsection{コマ主題細目}
\begin{description}
	\item[母集団と標本] 推測統計学の基本的な概念である母集団と標本の関係を理解し，統計学的な推定(点推定，区間推定)の位置付けについて学ぶ。\footnote{本書では検定をモデル比較の一種として扱うので，ここでは言及しない。}用語としての母数(母平均，母分散，母比率など)と標本統計量(標本平均，標本分散，標本比率など)，および標本統計量の実現値といった用語の意味するところも慎重に理解しなければならない。特に母集団分布，標本分布，サンプリング分布の違いについて理解することが重要である。
	      \begin{flushright}
		      $\to$ \textcite{Yamada2004}のP.68--73
	      \end{flushright}

	\item[母集団分布とモデル] ここでは母集団に正規分布を仮定する``モデル''について考える。統計的な推測は基本的に不良設定問題であり，何らかの前提をおかないと解くことができない問題がほとんどである。心理学の場合は母集団分布を正規分布とする仮定が選ばれることが多く，この場合は標本統計量の分布も計算しやすくなる。現実の分布が必ずしも正規分布に従わない場合でも，中心極限定理により標本平均は正規分布に近似できることが，この仮定の有用性を支えている。
	      \begin{flushright}
		      $\to$ \textcite{Yamada2004}のP.74--79.
	      \end{flushright}

	\item[正規分布の別の意味] 正規分布はここまで，誤差の分布として解説してきたが，複数の要因が積み重なった時の全体的傾向として理解することもできる。平均値が集団の代表値としての意味もあること，また中心極限定理によって分布によらずその標本平均値が正規分布に従うことが示される。中心極限定理は，母集団分布がどのような形であっても，標本サイズが十分大きければ標本平均の分布は正規分布に近づくという重要な性質であり，推測統計学の基礎となる定理である。
	      \begin{flushright}
		      中心極限定理については$\to$ \textcite{Minamoto2015}のP.180--181.
	      \end{flushright}

	\item[標本分布と標準誤差] 母集団が正規分布に従うとすると，標本統計量も正規分布に従うことが導出できる。ここでは標本統計量としての平均値を例にあげ，「標本統計量の分布」のことを標本分布と呼ぶこと，また標本分布の標準偏差をとくに標準誤差と呼ぶことを理解する。標準誤差は標本平均のばらつきを表す指標であり，サンプルサイズの平方根に反比例することを理解する。たとえばサンプルサイズが大きくなったら，小さくなったらどのような値が得られるのかについて，具体的な例をみながら考えると理解が進む。
	      \begin{flushright}
		      $\to$ 標本分布については\textcite{Yamada2004}のP.90--97.
	      \end{flushright}

	\item[標準正規分布と統計的推論] 標準正規分布は統計的推論において中心的な役割を果たす。標本統計量を標準化することで標準正規分布に対応させ，その確率的な性質を利用して推定や後の検定を行うことができる。標準正規分布表の使い方や，標準正規分布に基づく信頼区間の構成方法について学ぶ。また，標準正規分布を用いた統計的推論の基本的な考え方と，これが後の検定につながることを理解する。
	      \begin{flushright}
		      $\to$ \textcite{Yamada2004}のP.86--89，\textcite{Kawabata201408}のP.36--44.
	      \end{flushright}
\end{description}
\subsubsection{キーワード}
\begin{itemize}
	\item 母集団と標本
	\item 標本分布と標本統計量
	\item 中心極限定理
	\item 標準誤差
\end{itemize}
\subsection{授業情報}
\paragraph{コマの展開方法}
講義
\paragraph{標準シラバスにおける位置づけ}
科目番号5；心理学統計法；(2)統計に関する基礎的な知識；A 統計分析の基礎：母集団と標本
\subsubsection{予習・復習課題}
\paragraph{予習}
確率と確率分布の基本概念を復習しておくこと。特に標準正規分布の特徴とその利用方法について確認しておくと理解が深まる。また，ここでは\textcite{Yamada2004}のP.68-97が丁寧に解説されているので，事前に一読しておくことを強く勧める。
\paragraph{復習}
いよいよ記述統計をこえ，直接知り得ないものについて推論する段階に進む。基本的に解けない問題を解こうとするようなものであり，仮定やモデル，みなしている箇所について自覚的でなければならない。今回の講義の中で，何が仮定され，何が理論的に導出されたものなのかを区別できるかどうか，ノートを整理しておくと良いだろう。また推測統計学の基本的概念である，母集団と標本の関係と推定の手続きについては，\textcite{Yamada2004}のテキストが最も丁寧で適切な記述をしているので，是非該当箇所を通読し，概念や用語を間違えないように注意されたし。

以下の点について理解を深めておくことも有効である：
\begin{enumerate}
\item 標本サイズと標準誤差の関係
\item 中心極限定理の意味とその応用
\item 母集団分布を正規分布と仮定することの意義と限界
\item 点推定と区間推定の違いと用途
\end{enumerate}